% !TeX root = ../main.tex
\section{Introduction}
\label{sec:introduction}

Reservoir flow, reactive transport, and geomechanics are usually formulated
from different primary variables and with different levels of thermodynamic
resolution.  Compositional simulators evolve phase/component inventories and
close phase fluxes with Darcy laws, phase behavior, and capillary relations.
Poromechanics instead emphasizes deformation, pore pressure, and effective
stress, while electrochemical and geochemical models introduce charge,
diffusion, and reaction through additional constitutive laws.  These
specializations are effective in their intended regimes, but their direct
coupling can obscure whether a pressure, source, flux, or interaction force is
derived from a balance law, a constraint, or a dissipative closure.

Several established theories supply the ingredients needed for a more unified
description.  Reacting-mixture variational theories already provide
phase-attached motions, material conversion, constituent momentum production,
and Darcy-type porous-flow reductions
\cite{bedford1979variational,drumheller1980thermomechanical,bedford1983theories,drumheller2000}.
Macroscopic porous-medium thermodynamics derives multiphase balances,
capillary restrictions, and generalized transport laws from averaged energy
and entropy statements
\cite{hassanizadeh1993thermodynamic,gray1999thermodynamics}, while
poromechanics identifies effective-stress and equivalent-pressure structures
\cite{coussy2004poromechanics,kim2013rigorous}.  Seguin and Walkington give a
nonreactive multicomponent, multiphase poroelastic theory in a solid observer
\cite{seguin2019multicomponent}, and Foster and Xu derive a fixed-pressure
nonlinear Biot coefficient for a finite-deformation binary mixture
\cite{foster2025}.  Constrained reactive-solid theories permit composition and
constituent reference configurations to evolve during growth, plasticity, and
damage \cite{zimmerman2021constrainedreactive,ateshian2022constrainedreactive}.
For reacting and diffusing mixtures, force--flux restrictions follow from
continuum thermodynamics and Onsager stationarity
\cite{bothe2015continuum,doi2011onsager,wang2021onsager}; nonisothermal
chemical-potential gradients also require explicit attention to potential
gauges \cite{goldobin2016thermal}.  Finally, electroelastic potential theory
relates electric enthalpy, displacement, and Maxwell stress
\cite{mcmeeking2007virtual_work_electrostatic,dorfmann2017nonlinear_electroelasticity}.
Bennethum and Cushman derive the macroscale electroquasistatic field equations
for multicomponent, multiphase porous media in Part~I and the corresponding
constitutive restrictions for phase stresses, diffusion, and Darcy flow in
Part~II \cite{bennethum2002electroI,bennethum2002electroII}.

These precedents contain important parts of the present model, but the
inspected formulations do not present the six coupled extensions developed
below.  The novelty claimed here is therefore not the isolated use of
nonequilibrium thermodynamics, a Darcy law, electric stress, or reactive-solid
kinematics.  It is the integration of those structures in one open-system
phase/component derivation and the specific extensions that follow.  The
theory does not replace more detailed treatments of sharp interfaces,
pore-scale geometry, phase behavior, or material calibration.

The derivation has three steps with distinct roles.  A
constrained Hamilton--d'Alembert variation determines the reversible phase
coordinates, storage multipliers, momentum balances, and electrostatic field
equation.  A Coleman--Noll exploitation of two subsystem energy balances and
the entropy inequality determines the admissible stresses, pressures,
entropies, component potentials, and residual force--rate products.  An
Onsager stationarity principle then supplies one near-equilibrium closure for
reaction and component-relative transport.  Keeping these operations separate
prevents a balance constraint from being used as a constitutive law and makes
each retained dissipative mechanism visible.

Within that construction, the paper makes six specific contributions.
%
\begin{enumerate}
	\item The open-system momentum balance retains the momentum carried by mass
	      converted into each phase.  Its fluid reduction gives a generalized
	      Darcy law containing equivalent- and capillary-pressure gradients, body
	      force, Maxwell stress, inertia, and conversion-insertion momentum.  Net
	      phase conversion changes the algebraic resistance and mobility used to
	      eliminate the phase velocity; the exact nonsingularity and stronger
	      positive-resistance conditions are distinguished from dissipation of the
	      underlying drag.

	\item A fixed-equivalent-pressure Legendre transform extends the nonlinear
	      Biot construction from a binary mixture to multiple reacting solid
	      phases at finite deformation.  The phase coefficient is evaluated at the
	      current reacted state and changes through solid mass, composition,
	      distension, intrinsic specific volume, stress-free deformation, and
	      internal variables.  An additive component decomposition is obtained
	      only on the stated neutral, common-pressure component-EOS path; the
	      reacted phase coefficient is the general result.

	\item The thermal and electrical construction gives gauge-invariant
	      electrochemical transport and reaction forces when the mobile-fluid and
	      solid subsystems have different temperatures.  The
	      temperature-weighted neutral-potential gradient remains separate from
	      the electric field, while the usual electrochemical-potential gradient
	      is recovered in the isothermal limit.  The same charge-rate and field
	      power accounting prevents electrical body force or source power from
	      being counted twice.

	\item Capillarity and electrostatics are generated from the same constrained
	      stored-energy construction.  The saturation constraint produces
	      capillary-pressure differences, and the phase electrical enthalpies
	      contribute a definite difference term to those pressure jumps while
	      also generating electric displacement, Maxwell stress, and dielectric
	      pressure.  The resulting electrocapillary law and electrically corrected
	      equivalent pore pressure are therefore consistent with the phase
	      momentum equations, rather than appended as independent closures.

	\item Variation of the component material measures gives a generalized
	      composition-stationarity system.  At equilibrium it recovers equality of
	      the appropriate electrochemical potentials; away from equilibrium it
	      determines phase-reference transfer-work levels.  Consequently, the
	      conversion force for phase change, dissolution, precipitation, and pore
	      filling differs from the traditional stoichiometric affinity by an
	      explicit phase-transfer-work term.

	\item A single phase/component balance structure couples mobile phases to
	      multiple load-bearing solid phases that share the skeleton motion but
	      retain phase-specific true deformation, distension, stress-free maps,
	      composition, and internal variables.  It thereby represents
	      pore-filling conversion as accumulated skeleton-attached solid mass that
	      changes phase volume and can alter stress-free deformation, elasticity,
	      and plastic flow.  This extends reactive-solid kinematics to an open
	      multiphase porous medium with simultaneous transport and pressure
	      coupling.
\end{enumerate}
%
These features are intended to preserve a direct route to familiar engineering
models.  The theory is pulled back to the solid reference configuration and
then reduced explicitly to classical compositional flow, a black-oil model, a
Nernst--Planck--Darcy system, and finite- and small-deformation
single-solid/single-fluid Biot equations.  Each reduction states which
electrical, thermal, conversion, kinematic, or constitutive terms are removed.
The resulting correspondence is a consistency check on the general
formulation, not a claim that the general model supplies the equations of
state, relative permeabilities, kinetic data, or numerical algorithms required
by a particular application.

The scope is deliberately local and continuum-scale.  Constitutive functions
are assumed sufficiently smooth for the phases and components that are present,
with positive phase densities, volume fractions, and temperatures.  Separate
switching conditions are needed when a phase or component appears or
disappears.  The interfacial energy gives an algebraic local-equilibrium
capillary closure; dynamic capillary pressure, hysteresis, and independent
interfacial area are not included.  The electroquasistatic state describes an isotropic
dielectric response without permanent polarization, dielectric relaxation, or
electromagnetic-wave dynamics.  Source-carried energy and entropy use a
local-equilibrium insertion convention.  The general two-temperature reaction
force uses the neutral transfer coefficient so that charged transfer between
the thermal subsystems remains independent of the electric-potential gauge;
its physical reaction-energy exchange may be Onsager-coupled subject to the
stated mechanismwise dissipation condition.  The displayed quadratic Onsager
laws are near-equilibrium examples rather than unique kinetic models.  Boundary
data, gauges, and compatibility conditions are stated, but analysis of
well-posedness and numerical validation is deferred.

Section~\ref{sec:technical_setting} fixes the notation, observers, constitutive
scope, and data requirements.  The phase/component mass and charge balances
follow.  The Hamilton--d'Alembert--Onsager construction is then developed.
Section~\ref{sec:constitutive_models_entropy_inequality} introduces the solid
kinematics, performs the Coleman--Noll restriction, and derives the
dissipative closures, Biot transformations, and crystallization specialization.
The solid-reference equations are then collected for reservoir implementation,
and the final reductions establish correspondence with the classical models
listed above.
